%% kapa: exam=1 round=25 year=2014 period=2 form=A subject=accounting subject_ko=회계학 seq=23 number=63
%% source: Q-Net 감정평가사 기출문제 (cst00309) artlSeq=5101207 file=extracted/1008/5101207/2014년 제25회 감정평가사 1차시험 기출문제/2교시(A형).hwp
%% conversion: kordoc 4.24.51 -> tools/gen1.py (auto); review=proofread-sonnet; flags=none
\begin{question}{63}
(주)대한은 자사가 소유하고 있는 기계장치를 (주)세종이 소유하고 있는 차량운반구와 교환하였다. 두 기업의 유형자산에 관한 정보와 세부 거래 내용은 다음과 같다.
\begin{framedtext}
○ 이 교환은 상업적 실질이 있는 거래이다. \par
○ (주)대한의 기계장치 공정가치가 더 명백하다. \par
○ (주)세종은 (주)대한에게 공정가치의 차이인 ￦5,000을 지급하였다. \par
{\centering\small
\begin{tabular}{@{}l@{\hspace{1.5em}}r@{\hspace{1.5em}}r@{}}
\hline
 & (주)대한 & (주)세종 \\
 & \uline{기계장치} & \uline{차량운반구} \\
취득원가 & ￦50,000 & ￦50,000 \\
감가상각누계액 & 30,000 & 20,000 \\
공정가치 & 30,000 & 25,000 \\
현금지급액 & 0 & 5,000 \\
현금수취액 & 5,000 & 0 \\
\hline
\end{tabular}\par}
\end{framedtext}
이 거래와 관련한 설명 중 옳은 것은?
\begin{choices}
  \item (주)대한은 이 교환거래와 관련하여 유형자산처분이익 ￦5,000을 인식해야 한다.
  \item (주)대한이 새로 취득한 차량운반구의 취득원가는 ￦30,000이다.
  \item (주)세종은 이 교환거래와 관련하여 유형자산처분이익 ￦5,000을 인식해야 한다.
  \item (주)세종이 새로 취득한 기계장치의 취득원가는 ￦30,000이다.
  \item (주)대한과 (주)세종 모두 유형자산처분손익을 인식하지 않는다.
\end{choices}
\end{question}
\answer{4}
